% SPEC-TH-020 --- Guardrails: cost caps, depth limits, hallucination checks.
% Generic spec for the Thalamus cortex orchestrator's safety envelope.

% Shared preamble for all spec documents.
% Include from each spec: % Shared preamble for all spec documents.
% Include from each spec: % Shared preamble for all spec documents.
% Include from each spec: \input{../preamble.tex} (adjust depth).

\documentclass[11pt,a4paper]{article}

\usepackage[utf8]{inputenc}
\usepackage[T1]{fontenc}
\usepackage{lmodern}
\usepackage[a4paper,margin=2.3cm]{geometry}
\usepackage{parskip}
\usepackage{microtype}
\usepackage[dvipsnames]{xcolor}
\usepackage{enumitem}
\usepackage{listings}
\usepackage{tcolorbox}
\usepackage{titlesec}
\usepackage{fancyhdr}
\usepackage{array}
\usepackage{longtable}
\usepackage{booktabs}
\usepackage{amsmath}
\usepackage{amssymb}
\usepackage{hyperref}
\usepackage{cleveref}

\tcbuselibrary{skins,breakable}

\definecolor{specBlue}{RGB}{30,64,132}
\definecolor{specGray}{RGB}{90,90,90}
\definecolor{specAccent}{RGB}{198,40,40}
\definecolor{specGreen}{RGB}{30,110,60}

\hypersetup{
  colorlinks=true,
  linkcolor=specBlue,
  urlcolor=specBlue,
  citecolor=specBlue,
  pdfborder={0 0 0}
}

\titleformat{\section}{\Large\bfseries\color{specBlue}}{\thesection}{0.75em}{}
\titleformat{\subsection}{\large\bfseries\color{specBlue!80}}{\thesubsection}{0.6em}{}

\makeatletter
\newcommand{\specID}[1]{\def\@specID{#1}}
\newcommand{\specTitle}[1]{\def\@specTitle{#1}}
\newcommand{\specStatus}[1]{\def\@specStatus{#1}}
\newcommand{\specVersion}[1]{\def\@specVersion{#1}}
\newcommand{\specOwner}[1]{\def\@specOwner{#1}}
\newcommand{\specTraces}[1]{\def\@specTraces{#1}}

\specID{SPEC-XXX}
\specTitle{Untitled Spec}
\specStatus{DRAFT}
\specVersion{0.1}
\specOwner{Jeremie Nunez}
\specTraces{}

\newcommand{\makespecheader}{%
  \begin{center}
    {\LARGE\bfseries\color{specBlue}\@specTitle}\\[0.4em]
    {\small\color{specGray}\texttt{\@specID} \quad v\@specVersion \quad \textsc{\@specStatus} \quad \@specOwner}
  \end{center}
  \vspace{0.4em}
  \hrule
  \vspace{0.8em}
}

\pagestyle{fancy}
\fancyhf{}
\fancyhead[L]{\small\color{specGray}\@specID}
\fancyhead[R]{\small\color{specGray}\@specTitle}
\fancyfoot[C]{\small\color{specGray}\thepage}
\renewcommand{\headrulewidth}{0pt}
\makeatother

\newtcolorbox{invariant}{
  colback=specBlue!5, colframe=specBlue!75,
  title=Invariant, fonttitle=\bfseries, breakable,
  boxrule=0.5pt, arc=2pt
}

\newtcolorbox{scenario}[1]{
  colback=white, colframe=specBlue!50,
  title={Scenario: #1}, fonttitle=\bfseries, breakable,
  boxrule=0.5pt, arc=2pt
}

\newtcolorbox{ac}[1]{
  colback=specAccent!5, colframe=specAccent!60,
  title={Acceptance Criterion #1}, fonttitle=\bfseries, breakable,
  boxrule=0.5pt, arc=2pt
}

\newtcolorbox{nongoal}{
  colback=specGray!8, colframe=specGray!60,
  title=Non-Goal, fonttitle=\bfseries, breakable,
  boxrule=0.5pt, arc=2pt
}

\lstdefinestyle{codestyle}{
  basicstyle=\ttfamily\small,
  breaklines=true,
  showstringspaces=false,
  frame=single,
  framesep=4pt,
  rulecolor=\color{specBlue!30},
  backgroundcolor=\color{specBlue!2},
  xleftmargin=0.5em,
  xrightmargin=0.5em,
  keywordstyle=\color{specBlue}\bfseries,
  commentstyle=\color{specGray}\itshape,
  stringstyle=\color{specGreen}
}
\lstset{
  inputencoding=utf8,
  extendedchars=true,
  literate=
    {—}{{--}}1
    {–}{{-}}1
    {…}{{...}}1
    {’}{{'}}1
    {‘}{{'}}1
    {“}{{``}}1
    {”}{{''}}1
    {é}{{\'e}}1
    {è}{{\`e}}1
    {ê}{{\^e}}1
    {à}{{\`a}}1
    {â}{{\^a}}1
    {ç}{{\c{c}}}1
    {ô}{{\^o}}1
    {→}{{$\to$}}1
    {←}{{$\leftarrow$}}1
    {×}{{$\times$}}1
    {≥}{{$\geq$}}1
    {≤}{{$\leq$}}1
    {≠}{{$\neq$}}1
    {∈}{{$\in$}}1
    {⌈}{{$\lceil$}}1
    {⌉}{{$\rceil$}}1
    {∞}{{$\infty$}}1
    {α}{{$\alpha$}}1
    {β}{{$\beta$}}1
}
\lstdefinelanguage{TypeScript}{
  keywords={const,let,var,function,return,if,else,for,while,class,interface,type,extends,implements,import,export,from,as,async,await,new,this,true,false,null,undefined,void,any,unknown,never,string,number,boolean,readonly,private,public,protected,enum},
  sensitive=true,
  comment=[l]{//},
  morecomment=[s]{/*}{*/},
  morestring=[b]",
  morestring=[b]',
  morestring=[b]`
}
\lstset{style=codestyle}

\newcommand{\Given}{\textbf{\color{specBlue}Given} }
\newcommand{\When}{\textbf{\color{specBlue}When} }
\newcommand{\Then}{\textbf{\color{specBlue}Then} }
\providecommand{\And}{}
\renewcommand{\And}{\textbf{\color{specBlue}And} }
 (adjust depth).

\documentclass[11pt,a4paper]{article}

\usepackage[utf8]{inputenc}
\usepackage[T1]{fontenc}
\usepackage{lmodern}
\usepackage[a4paper,margin=2.3cm]{geometry}
\usepackage{parskip}
\usepackage{microtype}
\usepackage[dvipsnames]{xcolor}
\usepackage{enumitem}
\usepackage{listings}
\usepackage{tcolorbox}
\usepackage{titlesec}
\usepackage{fancyhdr}
\usepackage{array}
\usepackage{longtable}
\usepackage{booktabs}
\usepackage{amsmath}
\usepackage{amssymb}
\usepackage{hyperref}
\usepackage{cleveref}

\tcbuselibrary{skins,breakable}

\definecolor{specBlue}{RGB}{30,64,132}
\definecolor{specGray}{RGB}{90,90,90}
\definecolor{specAccent}{RGB}{198,40,40}
\definecolor{specGreen}{RGB}{30,110,60}

\hypersetup{
  colorlinks=true,
  linkcolor=specBlue,
  urlcolor=specBlue,
  citecolor=specBlue,
  pdfborder={0 0 0}
}

\titleformat{\section}{\Large\bfseries\color{specBlue}}{\thesection}{0.75em}{}
\titleformat{\subsection}{\large\bfseries\color{specBlue!80}}{\thesubsection}{0.6em}{}

\makeatletter
\newcommand{\specID}[1]{\def\@specID{#1}}
\newcommand{\specTitle}[1]{\def\@specTitle{#1}}
\newcommand{\specStatus}[1]{\def\@specStatus{#1}}
\newcommand{\specVersion}[1]{\def\@specVersion{#1}}
\newcommand{\specOwner}[1]{\def\@specOwner{#1}}
\newcommand{\specTraces}[1]{\def\@specTraces{#1}}

\specID{SPEC-XXX}
\specTitle{Untitled Spec}
\specStatus{DRAFT}
\specVersion{0.1}
\specOwner{Jeremie Nunez}
\specTraces{}

\newcommand{\makespecheader}{%
  \begin{center}
    {\LARGE\bfseries\color{specBlue}\@specTitle}\\[0.4em]
    {\small\color{specGray}\texttt{\@specID} \quad v\@specVersion \quad \textsc{\@specStatus} \quad \@specOwner}
  \end{center}
  \vspace{0.4em}
  \hrule
  \vspace{0.8em}
}

\pagestyle{fancy}
\fancyhf{}
\fancyhead[L]{\small\color{specGray}\@specID}
\fancyhead[R]{\small\color{specGray}\@specTitle}
\fancyfoot[C]{\small\color{specGray}\thepage}
\renewcommand{\headrulewidth}{0pt}
\makeatother

\newtcolorbox{invariant}{
  colback=specBlue!5, colframe=specBlue!75,
  title=Invariant, fonttitle=\bfseries, breakable,
  boxrule=0.5pt, arc=2pt
}

\newtcolorbox{scenario}[1]{
  colback=white, colframe=specBlue!50,
  title={Scenario: #1}, fonttitle=\bfseries, breakable,
  boxrule=0.5pt, arc=2pt
}

\newtcolorbox{ac}[1]{
  colback=specAccent!5, colframe=specAccent!60,
  title={Acceptance Criterion #1}, fonttitle=\bfseries, breakable,
  boxrule=0.5pt, arc=2pt
}

\newtcolorbox{nongoal}{
  colback=specGray!8, colframe=specGray!60,
  title=Non-Goal, fonttitle=\bfseries, breakable,
  boxrule=0.5pt, arc=2pt
}

\lstdefinestyle{codestyle}{
  basicstyle=\ttfamily\small,
  breaklines=true,
  showstringspaces=false,
  frame=single,
  framesep=4pt,
  rulecolor=\color{specBlue!30},
  backgroundcolor=\color{specBlue!2},
  xleftmargin=0.5em,
  xrightmargin=0.5em,
  keywordstyle=\color{specBlue}\bfseries,
  commentstyle=\color{specGray}\itshape,
  stringstyle=\color{specGreen}
}
\lstset{
  inputencoding=utf8,
  extendedchars=true,
  literate=
    {—}{{--}}1
    {–}{{-}}1
    {…}{{...}}1
    {’}{{'}}1
    {‘}{{'}}1
    {“}{{``}}1
    {”}{{''}}1
    {é}{{\'e}}1
    {è}{{\`e}}1
    {ê}{{\^e}}1
    {à}{{\`a}}1
    {â}{{\^a}}1
    {ç}{{\c{c}}}1
    {ô}{{\^o}}1
    {→}{{$\to$}}1
    {←}{{$\leftarrow$}}1
    {×}{{$\times$}}1
    {≥}{{$\geq$}}1
    {≤}{{$\leq$}}1
    {≠}{{$\neq$}}1
    {∈}{{$\in$}}1
    {⌈}{{$\lceil$}}1
    {⌉}{{$\rceil$}}1
    {∞}{{$\infty$}}1
    {α}{{$\alpha$}}1
    {β}{{$\beta$}}1
}
\lstdefinelanguage{TypeScript}{
  keywords={const,let,var,function,return,if,else,for,while,class,interface,type,extends,implements,import,export,from,as,async,await,new,this,true,false,null,undefined,void,any,unknown,never,string,number,boolean,readonly,private,public,protected,enum},
  sensitive=true,
  comment=[l]{//},
  morecomment=[s]{/*}{*/},
  morestring=[b]",
  morestring=[b]',
  morestring=[b]`
}
\lstset{style=codestyle}

\newcommand{\Given}{\textbf{\color{specBlue}Given} }
\newcommand{\When}{\textbf{\color{specBlue}When} }
\newcommand{\Then}{\textbf{\color{specBlue}Then} }
\providecommand{\And}{}
\renewcommand{\And}{\textbf{\color{specBlue}And} }
 (adjust depth).

\documentclass[11pt,a4paper]{article}

\usepackage[utf8]{inputenc}
\usepackage[T1]{fontenc}
\usepackage{lmodern}
\usepackage[a4paper,margin=2.3cm]{geometry}
\usepackage{parskip}
\usepackage{microtype}
\usepackage[dvipsnames]{xcolor}
\usepackage{enumitem}
\usepackage{listings}
\usepackage{tcolorbox}
\usepackage{titlesec}
\usepackage{fancyhdr}
\usepackage{array}
\usepackage{longtable}
\usepackage{booktabs}
\usepackage{amsmath}
\usepackage{amssymb}
\usepackage{hyperref}
\usepackage{cleveref}

\tcbuselibrary{skins,breakable}

\definecolor{specBlue}{RGB}{30,64,132}
\definecolor{specGray}{RGB}{90,90,90}
\definecolor{specAccent}{RGB}{198,40,40}
\definecolor{specGreen}{RGB}{30,110,60}

\hypersetup{
  colorlinks=true,
  linkcolor=specBlue,
  urlcolor=specBlue,
  citecolor=specBlue,
  pdfborder={0 0 0}
}

\titleformat{\section}{\Large\bfseries\color{specBlue}}{\thesection}{0.75em}{}
\titleformat{\subsection}{\large\bfseries\color{specBlue!80}}{\thesubsection}{0.6em}{}

\makeatletter
\newcommand{\specID}[1]{\def\@specID{#1}}
\newcommand{\specTitle}[1]{\def\@specTitle{#1}}
\newcommand{\specStatus}[1]{\def\@specStatus{#1}}
\newcommand{\specVersion}[1]{\def\@specVersion{#1}}
\newcommand{\specOwner}[1]{\def\@specOwner{#1}}
\newcommand{\specTraces}[1]{\def\@specTraces{#1}}

\specID{SPEC-XXX}
\specTitle{Untitled Spec}
\specStatus{DRAFT}
\specVersion{0.1}
\specOwner{Jeremie Nunez}
\specTraces{}

\newcommand{\makespecheader}{%
  \begin{center}
    {\LARGE\bfseries\color{specBlue}\@specTitle}\\[0.4em]
    {\small\color{specGray}\texttt{\@specID} \quad v\@specVersion \quad \textsc{\@specStatus} \quad \@specOwner}
  \end{center}
  \vspace{0.4em}
  \hrule
  \vspace{0.8em}
}

\pagestyle{fancy}
\fancyhf{}
\fancyhead[L]{\small\color{specGray}\@specID}
\fancyhead[R]{\small\color{specGray}\@specTitle}
\fancyfoot[C]{\small\color{specGray}\thepage}
\renewcommand{\headrulewidth}{0pt}
\makeatother

\newtcolorbox{invariant}{
  colback=specBlue!5, colframe=specBlue!75,
  title=Invariant, fonttitle=\bfseries, breakable,
  boxrule=0.5pt, arc=2pt
}

\newtcolorbox{scenario}[1]{
  colback=white, colframe=specBlue!50,
  title={Scenario: #1}, fonttitle=\bfseries, breakable,
  boxrule=0.5pt, arc=2pt
}

\newtcolorbox{ac}[1]{
  colback=specAccent!5, colframe=specAccent!60,
  title={Acceptance Criterion #1}, fonttitle=\bfseries, breakable,
  boxrule=0.5pt, arc=2pt
}

\newtcolorbox{nongoal}{
  colback=specGray!8, colframe=specGray!60,
  title=Non-Goal, fonttitle=\bfseries, breakable,
  boxrule=0.5pt, arc=2pt
}

\lstdefinestyle{codestyle}{
  basicstyle=\ttfamily\small,
  breaklines=true,
  showstringspaces=false,
  frame=single,
  framesep=4pt,
  rulecolor=\color{specBlue!30},
  backgroundcolor=\color{specBlue!2},
  xleftmargin=0.5em,
  xrightmargin=0.5em,
  keywordstyle=\color{specBlue}\bfseries,
  commentstyle=\color{specGray}\itshape,
  stringstyle=\color{specGreen}
}
\lstset{
  inputencoding=utf8,
  extendedchars=true,
  literate=
    {—}{{--}}1
    {–}{{-}}1
    {…}{{...}}1
    {’}{{'}}1
    {‘}{{'}}1
    {“}{{``}}1
    {”}{{''}}1
    {é}{{\'e}}1
    {è}{{\`e}}1
    {ê}{{\^e}}1
    {à}{{\`a}}1
    {â}{{\^a}}1
    {ç}{{\c{c}}}1
    {ô}{{\^o}}1
    {→}{{$\to$}}1
    {←}{{$\leftarrow$}}1
    {×}{{$\times$}}1
    {≥}{{$\geq$}}1
    {≤}{{$\leq$}}1
    {≠}{{$\neq$}}1
    {∈}{{$\in$}}1
    {⌈}{{$\lceil$}}1
    {⌉}{{$\rceil$}}1
    {∞}{{$\infty$}}1
    {α}{{$\alpha$}}1
    {β}{{$\beta$}}1
}
\lstdefinelanguage{TypeScript}{
  keywords={const,let,var,function,return,if,else,for,while,class,interface,type,extends,implements,import,export,from,as,async,await,new,this,true,false,null,undefined,void,any,unknown,never,string,number,boolean,readonly,private,public,protected,enum},
  sensitive=true,
  comment=[l]{//},
  morecomment=[s]{/*}{*/},
  morestring=[b]",
  morestring=[b]',
  morestring=[b]`
}
\lstset{style=codestyle}

\newcommand{\Given}{\textbf{\color{specBlue}Given} }
\newcommand{\When}{\textbf{\color{specBlue}When} }
\newcommand{\Then}{\textbf{\color{specBlue}Then} }
\providecommand{\And}{}
\renewcommand{\And}{\textbf{\color{specBlue}And} }


\specID{SPEC-TH-020}
\specTitle{Guardrails --- Cost Caps, Depth Limits, Hallucination Checks}
\specStatus{DRAFT}
\specVersion{0.1}
\specOwner{Jeremie Nunez}

\begin{document}
\makespecheader

\section{Context}
The Thalamus orchestrator drives recursive LLM-backed cortices that query internal data, external sources and other cortices. Without a hard safety envelope, such loops can burn unbounded tokens, recurse without progress, or surface fabricated claims.

This spec defines \textbf{Guardrails}: the non-bypassable control layer that bounds every cortex invocation and every research chain. Enforcement lives in the orchestrator, not in the cortex prompts --- prompts can drift, code cannot. This is a sovereignty-critical surface: it is where operators retain control of cost, depth, and factuality against non-deterministic LLM behaviour.

Guardrails sit between the cortex executor (\texttt{packages/thalamus/src/cortices/executor.ts}) and the cortex LLM transport. All cost, depth and validation decisions pass through it; none of them are delegated to the model.

\section{Goals}
\begin{itemize}[leftmargin=1.2em]
  \item Enforce a per-cortex \textbf{token and monetary budget} at call time, independent of what the cortex prompt asks for.
  \item Enforce a \textbf{depth limit} on any recursive or chained invocation, measured in iterations and in cumulative cost.
  \item Detect \textbf{hallucinations} via structured-output validation, schema conformance and source-citation requirements; reject or quarantine findings that fail.
  \item On breach, always surface a \textbf{partial result with an explicit breach marker}; never silently truncate.
  \item Emit \textbf{observability signals} (metrics, structured logs) on every breach and every quarantine, so operators can audit non-deterministic behaviour post-hoc.
\end{itemize}

\section{Non-Goals}
\begin{nongoal}
This spec does not define:
\begin{itemize}[leftmargin=1.2em]
  \item The content of cortex prompts (see SPEC-TH-011, cortex skills).
  \item Prompt-injection sanitation of inbound data payloads (see SPEC-TH-021, input sanitiser). Guardrails here govern \emph{orchestrator behaviour}, not \emph{payload cleanup}.
  \item Rate limiting at the HTTP edge (see SPEC-API-030).
  \item Long-term cost accounting across tenants (see SPEC-BILL-010).
  \item Model selection, fallback strategy, or retry policy at the transport layer (see SPEC-TH-014).
\end{itemize}
\end{nongoal}

\section{Invariants}

\begin{invariant}
\textbf{INV-1 --- Budget is non-bypassable.} Every cortex invocation is wrapped by \texttt{withGuardrails(budget, fn)}. It is impossible for a cortex to reach the LLM transport without a resolved, typed \texttt{Budget} object. Calling the transport directly is a lint-enforced violation.
\end{invariant}

\begin{invariant}
\textbf{INV-2 --- Budget is monotonic.} Cumulative token and monetary cost on a chain only ever increase; there is no reset, discount or retroactive rebate across iterations of the same research chain. Reported cost $\geq$ actual cost for any prefix of the chain.
\end{invariant}

\begin{invariant}
\textbf{INV-3 --- Depth is bounded by construction.} A research chain terminates in at most \texttt{maxIterationsPerChain} cycles, regardless of the novelty signal, LLM confidence, or cortex opinion. The orchestrator decrements the remaining-depth counter before dispatching each iteration; depth 0 short-circuits to a partial result.
\end{invariant}

\begin{invariant}
\textbf{INV-4 --- Breach is observable.} On every breach (budget, depth, validation), the result carries a \texttt{breach} field enumerating the triggered guard(s), the measured value and the threshold. A silent success path on breach is a defect.
\end{invariant}

\begin{invariant}
\textbf{INV-5 --- Unverifiable findings do not reach consumers.} A finding without a valid schema shape, or without at least one non-empty \texttt{evidence[]} citation pointing to an allow-listed source, is quarantined with \texttt{findingType = rejected\_unsourced} and never emitted to downstream stores.
\end{invariant}

\section{Interface}

The orchestrator exposes a single entry point that all cortex calls must flow through. The contract is total: every outcome maps to a typed result, including breaches.

\begin{lstlisting}[language=TypeScript]
// Budget resolved per complexity tier (simple | moderate | deep)
export interface Budget {
  maxIterations: number;   // hard cap on chain depth
  maxCost: number;         // monetary cap per chain, USD
  maxTokens: number;       // aggregate prompt+completion tokens
  maxWallClockMs: number;  // chain timeout
  confidenceTarget: number;
  coverageTarget: number;
}

export type BreachKind =
  | "budget_cost"
  | "budget_tokens"
  | "depth_iterations"
  | "wall_clock"
  | "schema_invalid"
  | "unsourced_finding"
  | "source_not_allowlisted";

export interface Breach {
  kind: BreachKind;
  measured: number | string;
  threshold: number | string;
  iteration: number;
  cortex: string;
}

export interface GuardedOutput<T> {
  result: T;                 // partial if breach.length > 0
  partial: boolean;
  breaches: Breach[];
  usage: { tokens: number; costUsd: number; iterations: number; };
}

export function withGuardrails<T>(
  budget: Budget,
  ctx: ChainContext,
  fn: (remaining: Budget) => Promise<T>,
): Promise<GuardedOutput<T>>;

export function validateFinding(
  raw: unknown,
  allowlist: SourceAllowlist,
): Finding | RejectedFinding;
\end{lstlisting}

Emitted metrics (Prometheus-compatible):
\begin{itemize}[leftmargin=1.2em]
  \item \texttt{thalamus\_guardrail\_breach\_total\{kind,cortex\}} --- counter
  \item \texttt{thalamus\_chain\_cost\_usd\{cortex,tier\}} --- histogram
  \item \texttt{thalamus\_chain\_iterations\{cortex,tier\}} --- histogram
  \item \texttt{thalamus\_finding\_rejected\_total\{reason,cortex\}} --- counter
\end{itemize}

\section{Scenarios}

\begin{scenario}{S-1 --- Nominal within budget}
\Given a chain with \texttt{budget = simple} (\$0.03, 2 iter) \\
\When iteration~1 costs \$0.008, produces 3 valid findings, and the novelty signal says stop \\
\Then the chain exits normally with \texttt{partial=false}, \texttt{breaches=[]}, and \texttt{usage.costUsd=0.008}.
\end{scenario}

\begin{scenario}{S-2 --- Monetary cap exceeded mid-chain}
\Given \texttt{budget.maxCost = 0.06} (moderate tier) \\
\And iterations 1--3 accumulated \$0.055 \\
\When iteration~4 is about to dispatch and the transport's pre-flight cost estimate is \$0.012 \\
\Then the orchestrator refuses to dispatch iteration~4, returns \texttt{GuardedOutput} with \texttt{partial=true}, \texttt{breaches=[\{kind:"budget\_cost", measured:0.055, threshold:0.06\}]}, and includes all findings accumulated through iteration~3. \\
\And metric \texttt{thalamus\_guardrail\_breach\_total\{kind="budget\_cost"\}} is incremented once.
\end{scenario}

\begin{scenario}{S-3 --- Token cap exceeded during iteration}
\Given a running iteration with \texttt{maxTokens=60\,000} remaining \\
\When the streaming LLM response crosses 60\,000 cumulative tokens mid-stream \\
\Then the transport stream is cancelled at the next token boundary \\
\And a partial finding set is parsed from the already-received prefix; any incomplete JSON object is dropped (not reconstructed) \\
\And the result carries \texttt{breaches=[\{kind:"budget\_tokens",\dots\}]}.
\end{scenario}

\begin{scenario}{S-4 --- Depth limit reached}
\Given \texttt{maxIterationsPerChain=8} and the chain at iteration~8 with an active novelty signal \\
\When the strategist cortex requests a 9th iteration \\
\Then the request is denied regardless of remaining cost budget \\
\And the output is marked \texttt{partial=true} with \texttt{kind:"depth\_iterations"}.
\end{scenario}

\begin{scenario}{S-5 --- Hallucination: finding without evidence}
\Given a cortex returns a finding with \texttt{evidence=[]} \\
\When \texttt{validateFinding} runs \\
\Then the finding is rejected, replaced by a \texttt{RejectedFinding} stub carrying \texttt{reason="unsourced"} \\
\And it is \emph{not} persisted to the knowledge graph \\
\And metric \texttt{thalamus\_finding\_rejected\_total\{reason="unsourced"\}} is incremented.
\end{scenario}

\begin{scenario}{S-6 --- Hallucination: non-allowlisted source}
\Given a finding cites a URL whose host is not in the per-cortex source allowlist \\
\When validation runs \\
\Then the finding is rejected with \texttt{reason="source\_not\_allowlisted"} and the offending URL is logged at \texttt{warn}. \\
\And the remainder of the finding set for that iteration is kept.
\end{scenario}

\begin{scenario}{S-7 --- Schema-invalid structured output}
\Given the LLM returns JSON that fails the cortex output Zod schema (missing \texttt{confidence} field) \\
\When the validator runs after one retry with schema-repair prompt \\
\Then on second failure, the iteration returns zero findings with \texttt{breaches=[\{kind:"schema\_invalid"\}]} \\
\And cost of the failed iteration is still charged to the chain budget (INV-2).
\end{scenario}

\begin{scenario}{S-8 --- Wall-clock timeout on a long iteration}
\Given a chain-level \texttt{maxWallClockMs=300\,000} \\
\When the current iteration's transport call exceeds 90\,000\,ms (\texttt{cortex.timeoutMs}) or the aggregate exceeds 300\,000\,ms \\
\Then the in-flight call is aborted via \texttt{AbortSignal}; partial findings from prior iterations are returned with \texttt{kind:"wall\_clock"}.
\end{scenario}

\begin{scenario}{S-9 --- Consecutive zero-finding iterations}
\Given \texttt{consecutiveZeroStop=2} \\
\When iteration~$n$ and iteration~$n{+}1$ both produce zero valid findings after validation \\
\Then the chain terminates with \texttt{partial=false} (not a breach --- a normal stop condition) \\
\And a structured log entry tags the stop reason as \texttt{"diminishing\_returns"}. This is \emph{not} a breach.
\end{scenario}

\section{Acceptance Criteria}

\begin{ac}{AC-1}
Given a chain where cumulative cost crosses \texttt{budget.maxCost} after iteration~$k$, the orchestrator MUST NOT dispatch iteration~$k{+}1$. The returned \texttt{GuardedOutput.partial} MUST be \texttt{true} and \texttt{breaches} MUST contain exactly one entry with \texttt{kind="budget\_cost"}.
\end{ac}

\begin{ac}{AC-2}
Given a chain at iteration equal to \texttt{budget.maxIterations}, any further dispatch MUST be refused with a \texttt{depth\_iterations} breach. This holds even if remaining cost budget is non-zero.
\end{ac}

\begin{ac}{AC-3}
Given a finding with \texttt{evidence.length === 0}, \texttt{validateFinding} MUST return a \texttt{RejectedFinding} and the finding MUST NOT appear in the persisted knowledge graph.
\end{ac}

\begin{ac}{AC-4}
Given a finding whose \texttt{evidence[i].url} host is not in the cortex's source allowlist, \texttt{validateFinding} MUST return a \texttt{RejectedFinding} with \texttt{reason="source\_not\_allowlisted"}.
\end{ac}

\begin{ac}{AC-5}
Given an LLM response that fails schema validation twice in a row for the same iteration, that iteration MUST yield zero findings and MUST still increment \texttt{usage.costUsd} with the attempted cost (monotonicity, INV-2).
\end{ac}

\begin{ac}{AC-6}
On every breach, exactly one increment of \texttt{thalamus\_guardrail\_breach\_total} with the matching \texttt{kind} and \texttt{cortex} labels MUST be observed. No breach MUST be silent; a unit test asserting breach-without-metric MUST fail the build.
\end{ac}

\begin{ac}{AC-7}
\texttt{GuardedOutput} MUST always be returned, even on breach. The orchestrator MUST NOT \texttt{throw} out of \texttt{withGuardrails} for any of the defined \texttt{BreachKind}s; thrown exceptions are reserved for infrastructure errors (DB unreachable, transport crash).
\end{ac}

\begin{ac}{AC-8}
Given a wall-clock overrun, the in-flight transport call MUST receive an \texttt{AbortSignal.abort()} within 500\,ms of the deadline, verified by an integration test that measures dispatch-to-abort latency.
\end{ac}

\begin{ac}{AC-9}
Partial results returned on breach MUST be structurally identical (same TypeScript type, same JSON shape) to successful results; consumers MUST be able to rely on \texttt{result} being safe to read regardless of \texttt{partial}'s value.
\end{ac}

\begin{ac}{AC-10}
The per-tier \texttt{Budget} (\texttt{simple}, \texttt{moderate}, \texttt{deep}) MUST be the \emph{only} source of numeric caps at runtime. No literal cost, iteration or token number MUST appear inside cortex prompts or cortex skill code --- enforced by a repository-level grep test.
\end{ac}

\section{Traceability}

\begin{center}
\begin{tabular}{@{}lll@{}}
\toprule
AC & Test file & Test name \\
\midrule
AC-1  & \texttt{packages/thalamus/test/guardrails.cost.spec.ts}      & \texttt{aborts chain on cost cap} \\
AC-2  & \texttt{packages/thalamus/test/guardrails.depth.spec.ts}     & \texttt{refuses dispatch at max depth} \\
AC-3  & \texttt{packages/thalamus/test/guardrails.validate.spec.ts}  & \texttt{rejects unsourced finding} \\
AC-4  & \texttt{packages/thalamus/test/guardrails.validate.spec.ts}  & \texttt{rejects non-allowlisted source} \\
AC-5  & \texttt{packages/thalamus/test/guardrails.validate.spec.ts}  & \texttt{charges cost on schema-invalid} \\
AC-6  & \texttt{packages/thalamus/test/guardrails.metrics.spec.ts}   & \texttt{emits one breach metric per breach} \\
AC-7  & \texttt{packages/thalamus/test/guardrails.contract.spec.ts}  & \texttt{never throws on defined breach kinds} \\
AC-8  & \texttt{packages/thalamus/test/guardrails.timeout.spec.ts}   & \texttt{aborts in-flight call within 500ms} \\
AC-9  & \texttt{packages/thalamus/test/guardrails.contract.spec.ts}  & \texttt{partial result has same shape} \\
AC-10 & \texttt{packages/thalamus/test/guardrails.nomagic.spec.ts}   & \texttt{no literal caps outside config} \\
\bottomrule
\end{tabular}
\end{center}

\section{Open Questions}
\begin{itemize}[leftmargin=1.2em]
  \item Should the cost cap be checked \emph{before} dispatch using a pre-flight token estimate, or only \emph{after} dispatch on reported usage? Current decision: pre-flight estimate + post-flight reconciliation. Needs validation against transport accuracy.
  \item Per-tenant vs.\ per-chain monetary caps: in scope for SPEC-BILL-010, out of scope here --- but the interface must not preclude tenant-level wrapping.
  \item Allowlist governance: who owns the per-cortex source allowlist? Proposed: declared in the cortex skill header, reviewed at skill-registration time. To be confirmed.
\end{itemize}

\end{document}
